\section{Methods} \label{moreno:sec:methods}

TODO methods.

\subsection{Software and Data Availability} \label{moreno:sec:software}

Simulation code, configuration files, and batch scripts for this work are available via Zenodo at \url{TODO}.
Postprocessing scripts and executable notebooks for this work are hosted separately at \url{TODO}.
Data and supplemental materials are available via the Open Science Framework \url{TODO} \cite{foster2017open}.

This project benefited significantly from open-source scientific software \cite{2020SciPy-NMeth,harris2020array,reback2020pandas,mckinney-proc-scipy-2010,waskom2021seaborn,hunter2007matplotlib,moreno2023teeplot,yang2025downstream,moreno2021signalgp,vostinar2024empirical,moreno2022hstrat,moreno2022best}.

\subsection{Wafer-Scale Engine}

The 3rd generation Cerebras WSE is a 755×1,170 mesh of 880,000 networked compute cores (Processing Elements or PEs; \cref{moreno:fig:wse}). PEs execute independently, and use specialized on-chip routing to stream 32-bit messages between neighbors with single-cycle latency. Each PE has 48kb of private on-chip memory with fast single-cycle read/write access \cite{buitrago2021neocortex}. PEs support standard CPU-like control flow, in addition to vectorized integer and floating-point operations.
In addition to deep learning frameworks (e.g., TensorFlow), Cerebras provides a full-featured Software Development Kit (SDK) \cite{selig2022cerebras} for general-purpose HPC \cite{jacquelin2022scalable,rocki2020fast}. The WSE device is programmed on a per-PE basis using the Cerebras Software Language (CSL). CSL organizes code within an event-driven framework, with the programmer defining tasks to be triggered in response to signals exchanged between—and within—PEs. Task handling occurs via hardware-level microthreading that supports cooperative concurrency.

As with other next-generation AI/ML accelerator platforms, work on WSE must account for scarce on-device memory (48kb per core), extreme decentralization (direct communication only among local neighbors), and restricted host-device connectivity (only at grid perimeter). Thus, translating HPC to these platforms requires software design directly informed by hardware.

\subsection{Hereditary Stratigraphy}

\cite{singhvi2025scalable,moreno2024structured,moreno2025testing}

\subsection{Benchmarking Experiments} \label{moreno:sec:methods-code-data-reproducibility-benchmarking-experiments}

Benchmarking experiments were performed on Michigan State University's High Performance Computing Center, a cluster of hundreds of heterogeneous x86 nodes linked with InfiniBand interconnects.
For multithread experiments, benchmarks for each thread count were collected from the same node.
For multiprocess experiments, each process was assigned to a distinct node in order to ensure results were representative of performance in a distributed context.
All multiprocess benchmarks were recorded from the same collection of nodes.
Hostnames are recorded for each benchmark data point.
For an exact accounting of hardware architectures used, these hostnames can be crossreferenced with a table included with the data that summarizes the cluster's node configurations.

Code for the distributed graph coloring benchmark is available at \url{https://github.com/mmore500/conduit} under \\ \texttt{demos/channel\_selection}.
Code for the digital evolution simulation benchmark is available at \url{https://github.com/mmore500/dishtiny}.
Exact versions of software used are recorded with each benchmark data point.
Data is available via the Open Science Framework at \url{https://osf.io/7jkgp/} \citep{foster2017open}.
Notebook code for all reported statistics and data visualizations is available at \url{https://hopth.ru/ci}.

\subsection{Quality of Service Experiments}

Quality of service experiments were carried out on Michigan State University's High Performance Computing Center lac cluster, consisting of 28-core Intel(R) Xeon(R) CPU E5-2680 v4 \@ 2.40GHz nodes.
All statistical comparisons are performed between observations from the same job allocation.
(Except in the case where intranode and internode configurations were compared, where experiments were performed on separate allocations using comparable nodes on the same cluster.)

Benchmarking experiments described in Section \ref{moreno:sec:methods-code-data-reproducibility-benchmarking-experiments} used a send/receive buffer size of 2.
However, due to the high communication intensity of the graph coloring problem with just one simulation element per CPU, quality of service experiments required a larger buffer size of 64 to maintain runtime stability.
In early work developing the Conduit library, we discovered that real-time messaging channels can enter a destabilizing positive feedback spiral when incoming messages take longer to handle (e.g., skip past or read) than sending messages.
Under such conditions, when a process exchanging messages from a partner process experiences a delay it sends fewer messages to that partner process.
Due to fewer incoming messages, the partner process can update more rapidly, increasing incoming message load on the delayed process.
This effect can snowball the partnership intended for even, two-way message exchange into effectively a unilateral producer-consumer relationship where (potentially unbounded) work piles up on the consumer.
To interrupt such a scenario, we use the bulk message pull call \verb|MPI_Testsome| to ensure fast message consumption under backlogged conditions.
So, receiver workload remains closer to constant under high traffic situations (instead of having to pull messages down one-by-one).
Larger receive buffer size, as configured for the quality of service experiments, increases the effectiveness of the bulk message consumption countermeasure.

Code for the distributed graph coloring benchmark is available at \url{https://github.com/mmore500/conduit} under \\ \texttt{demos/channel\_selection}.
Exact versions of software used are recorded with each benchmark data point.
Data is available via the Open Science Framework at \url{https://osf.io/72k5n/} \citep{foster2017open}.
Notebook code for all reported statistics is available at \url{https://hopth.ru/cj}.
